\documentclass[14pt,a4paper,oneside]{extbook}
\usepackage[utf8]{vietnam}
\usepackage{amsmath,amssymb,mathrsfs,fancyhdr,enumerate,multirow,makecell}
\usepackage{tkz-euclide,tikz-3dplot,tikz,tkz-tab}
\usetikzlibrary{arrows.meta,calc,intersections,angles,quotes,backgrounds,shapes.geometric,patterns,through,shadows,positioning}
\usepackage{yhmath}
\usepackage[top=0.7cm, bottom=0.7cm, left=1.0cm, right=1.0cm, headheight=16pt, headsep=3mm, footskip=5mm]{geometry}
\usepackage[hidelinks,unicode]{hyperref}
\usepackage{qrcode}
\usepackage[many]{tcolorbox}
\usepackage{fontawesome}

\pagestyle{fancy}
\fancyhf{}
\fancyhead[L]{\footnotesize\bfseries\color{teal!85!black}\faBook~SỔ TAY KIẾN THỨC TOÁN 4}
\fancyhead[R]{\footnotesize\bfseries\color{blue!85!black}Lớp Toán Thầy Lĩnh -- loptoanthaylinh.vn}
\fancyfoot[C]{\footnotesize\bfseries\color{gray!80!black}-- Sổ tay kiến thức: Bài 30. Thực hành vẽ hai đường thẳng song song --}
\renewcommand{\headrulewidth}{0.4pt}

\begin{document}

% --- BANNER TIÊU ĐỀ BÀI HỌC ---
\noindent
\begin{tcolorbox}[
  enhanced,
  colback=teal!8!white,
  colframe=teal!80!black,
  arc=3mm, boxrule=1pt,
  left=3mm, right=3mm, top=2mm, bottom=2mm,
  halign=center
]
  {\Large\bfseries\color{teal!90!black}\faBookmark\ BÀI 30: THỰC HÀNH VẼ HAI ĐƯỜNG THẲNG SONG SONG}\\[1.5pt]
  {\footnotesize\itshape\color{black!75} Quy tắc 2 lần vuông góc tạo nên song song, thao tác thước kẻ và ê-ke vẽ hình chữ nhật trên 1 trang A4}
\end{tcolorbox}

\vspace{2mm}

% --- KHỐI 1 & 2: LÝ THUYẾT & KĨ NĂNG ---
\noindent
\begin{minipage}[t]{0.49\textwidth}
  \begin{tcolorbox}[
    enhanced,
    title={\footnotesize\bfseries\color{white}\faPencil\ 1. Quy trình vẽ đường thẳng song song},
    colback=blue!2!white,
    colframe=blue!75!black,
    arc=2.5mm, boxrule=0.8pt, top=1.8mm, bottom=1.8mm, left=2mm, right=2mm
  ]
    \footnotesize
    Vẽ đường thẳng $CD$ đi qua điểm $M$ và song song với đường thẳng $AB$ cho trước:\\
    $\bullet$ \textbf{Bước 1:} Vẽ đường thẳng phụ $MN$ đi qua $M$ và \textbf{vuông góc} với $AB$ tại $N$.\\
    $\bullet$ \textbf{Bước 2:} Vẽ đường thẳng $CD$ đi qua $M$ và \textbf{vuông góc} với $MN$.\\
    $\bullet$ \textbf{Kết luận}: Vì $CD \perp MN$ và $AB \perp MN$ nên $CD \parallel AB$.
  \end{tcolorbox}
\end{minipage}%
\hfill
\begin{minipage}[t]{0.49\textwidth}
  \begin{tcolorbox}[
    enhanced,
    title={\footnotesize\bfseries\color{white}\faThLarge\ 2. Ứng dụng: Vẽ hình chữ nhật},
    colback=teal!2!white,
    colframe=teal!75!black,
    arc=2.5mm, boxrule=0.8pt, top=1.8mm, bottom=1.8mm, left=2mm, right=2mm
  ]
    \footnotesize
    Vẽ hình chữ nhật $ABCD$ có chiều dài $5\text{ cm}$, rộng $3\text{ cm}$:\\
    $\bullet$ \textbf{Bước 1:} Vẽ đoạn thẳng $AB = 5\text{ cm}$.\\
    $\bullet$ \textbf{Bước 2:} Vẽ đường thẳng vuông góc với $AB$ tại $A$, trên đó lấy $AD = 3\text{ cm}$.\\
    $\bullet$ \textbf{Bước 3:} Vẽ đường thẳng vuông góc với $AB$ tại $B$, trên đó lấy $BC = 3\text{ cm}$.\\
    $\bullet$ \textbf{Bước 4:} Nối $C$ với $D \rightarrow$ được hình chữ nhật $ABCD$.
  \end{tcolorbox}
\end{minipage}

\vspace{2mm}

% --- KHỐI 3: SƠ ĐỒ TƯ DUY TIKZ ---
\noindent
\begin{tcolorbox}[
  enhanced,
  title={\footnotesize\bfseries\color{white}\faSitemap\ SƠ ĐỒ TƯ DUY: QUY TẮC "2 LẦN VUÔNG GÓC TẠO NÊN 1 LẦN SONG SONG"},
  colback=yellow!2!white,
  colframe=orange!85!black,
  arc=2.5mm, boxrule=0.8pt, top=1.5mm, bottom=1.5mm, left=1mm, right=1mm,
  halign=center
]
  \begin{tikzpicture}[
    scale=0.92,
    stepb/.style={rectangle, rounded corners=2mm, draw=teal!80!black, fill=teal!10, thick, inner sep=2pt, font=\scriptsize\bfseries, align=center},
    arr/.style={-{Stealth[scale=0.85]}, thick, draw=black!70, shorten >= 2.5pt, shorten <= 2.5pt}
  ]
    % Hình vẽ thao tác
    \draw[thick, blue!85!black] (0.5, -0.8) -- (5.0, -0.8) node[right, font=\scriptsize\bfseries] {$AB$};
    \filldraw[black] (2.2, 0.8) circle (1.5pt) node[above, font=\scriptsize\bfseries] {$M$};

    % Đường phụ MN
    \draw[thick, dashed, orange!85!black] (2.2, 0.8) -- (2.2, -0.8) node[below, font=\scriptsize\bfseries] {$N$};
    \draw[thick, red!80!black] (2.2, -0.55) -- (2.45, -0.55) -- (2.45, -0.8);

    % Đường song song CD
    \draw[very thick, teal!85!black] (0.5, 0.8) -- (5.0, 0.8) node[right, font=\scriptsize\bfseries] {$CD$};
    \draw[thick, red!80!black] (2.2, 0.55) -- (2.45, 0.55) -- (2.45, 0.8);

    % Các bước bên phải
    \node[stepb, draw=blue!75!black, fill=blue!8, text width=9.5cm] (st1) at (10.5, 1.35) {\textbf{BƯỚC 1: KẺ ĐƯỜNG PHỤ VUÔNG GÓC}\\Dùng ê-ke vẽ đường thẳng $MN$ qua điểm $M$ và vuông góc với $AB$ tại $N$};
    \node[stepb, draw=teal!80!black, fill=teal!8, text width=9.5cm] (st2) at (10.5, 0.0) {\textbf{BƯỚC 2: KẺ ĐƯỜNG THẲNG VUÔNG GÓC THỨ HAI}\\Dùng ê-ke vẽ đường thẳng $CD$ qua điểm $M$ và vuông góc với đường phụ $MN$};
    \node[stepb, draw=purple!80!black, fill=purple!10, text width=9.5cm] (st3) at (10.5, -1.35) {\textbf{KẾT LUẬN: ĐƯỢC 2 ĐƯỜNG THẲNG SONG SONG}\\Đường thẳng $CD$ đi qua $M$ và song song tuyệt đối với $AB$ ($CD \parallel AB$)};

    \draw[arr] (st1) -- (st2);
    \draw[arr] (st2) -- (st3);
  \end{tikzpicture}
\end{tcolorbox}

\vspace{2mm}

% --- KHỐI 4: VÍ DỤ MINH HỌA & BẪY TRẮC NGHIỆM ---
\noindent
\begin{minipage}[t]{0.53\textwidth}
  \begin{tcolorbox}[
    enhanced,
    title={\footnotesize\bfseries\color{white}\faCheckCircle\ Ví dụ thực hành chuẩn mực},
    colback=green!2!white,
    colframe=green!60!black,
    arc=2.5mm, boxrule=0.8pt, top=1.8mm, bottom=1.8mm, left=2.2mm, right=2.2mm
  ]
    \footnotesize
    Vẽ hình vuông $MNPQ$ cạnh $4\text{ cm}$:\\
    $\bullet$ \textbf{B1}: Vẽ đoạn $MN = 4\text{ cm}$.\\
    $\bullet$ \textbf{B2}: Vẽ đoạn $MQ \perp MN$ tại $M$, $MQ = 4\text{ cm}$.\\
    $\bullet$ \textbf{B3}: Vẽ đoạn $NP \perp MN$ tại $N$, $NP = 4\text{ cm}$.\\
    $\bullet$ \textbf{B4}: Nối $P$ với $Q \rightarrow$ Được hình vuông $MNPQ$ có $PQ \parallel MN$ và $MQ \parallel NP$.
  \end{tcolorbox}
\end{minipage}%
\hfill
\begin{minipage}[t]{0.45\textwidth}
  \begin{tcolorbox}[
    enhanced,
    title={\footnotesize\bfseries\color{white}\faExclamationTriangle\ Mẹo tránh bẫy trắc nghiệm},
    colback=red!2!white,
    colframe=red!70!black,
    arc=2.5mm, boxrule=0.8pt, top=1.8mm, bottom=1.8mm, left=2.2mm, right=2.2mm
  ]
    \footnotesize
    $\bullet$ \textbf{Bẫy vẽ song song bằng ước lượng}:\\
    Không được dùng thước thẳng đặt "song song theo cảm tính", phải luôn qua thao tác trung gian \textbf{vuông góc} bằng ê-ke để đảm bảo chính xác.\\
    $\bullet$ Khi kẻ đường thẳng qua $M$, chú ý giữ chặt thước kẻ để ê-ke không bị trượt lệch.
  \end{tcolorbox}
\end{minipage}

\end{document}
